% =====================================================================
%  Chapter 29 review -- the 14-station revision contract + the Summit.
%  Revision-only for stations 1-11 and 13-14: tables, one-liners, pointers back
%  via \secref/\thmref/\lemref/\defref/\figref/\coderef/\probref.
%  Summit: CF 835F "Roads in the Kingdom" (2500, L5) -- a connected weighted
%  unicyclic graph (n vertices, n edges): close exactly one road so the kingdom
%  stays connected (i.e. delete a CYCLE edge) and MINIMISE the resulting tree's
%  diameter ("inconvenience").  Pipeline: peel leaves to find the unique cycle;
%  per cycle vertex the attached in-tree gives a height h_i and internal diameter
%  dtree_i (Ch.12 two-sweep / Ch.13 fold as black boxes); break the cycle at each
%  cycle edge -> a path over the k cycle vertices; the path's through-diameter is
%  max_{p<q} h_p + h_q + (S_q - S_p), answered per length-k window on a doubled
%  prefix array with a segment-tree monoid storing (best pair, max h+S, max h-S).
%  answer = max(max_i dtree_i, min_over_cuts windowDiameter).  Official sample
%  hand-replayed -> 5.  Reuses THREE chapter tools (cycle structure, contraction,
%  the DP) and TWO earlier-chapter tools (Ch.12 diameter/height, Ch.19/20 segment
%  tree).  Exports prob:c31-summit, lst:c31-summit, fig:c31-summit.
%  Mixed set EXACTLY 8, all live-verified this session against the CF problemset
%  API (titles/ratings/tags re-pulled) and absent from build/used_cf.txt; each
%  hint routes via the graph-to-tree signature (def:c31-signature) or REJECTS the
%  decomposition and names the true owner.  Two former mixed entries (CF 702E,
%  CF 231E) were PROMOTED into full in-section solved labs (prob:c31-funcpath in
%  29.3, prob:c31-cactuscount in 29.4), so their \practice entries were deleted
%  here to keep every CF URL appearing exactly once in the book.
% =====================================================================

\chaptersection{Chapter 29 Summary}

One sentence: \emph{a general graph is a tree plus a set of redundant edges, and this
chapter's entire craft is choosing which redundancy to delete -- the unavoidable
\textbf{edges} (bridges), the hinge \textbf{vertices} (articulation points), the
self-referential \textbf{arrows} of a function, or the \textbf{cycles} of a cactus --
so that what remains is a tree you can already run Ch.~12 and Ch.~13 on
\(\bigl(\secref{c31-bridges}, \secref{c31-block}, \secref{c31-functional},
\secref{c31-cactus}, \defref{c31-signature}\bigr)\).}  Chapter~28 ended by assuming the
input \emph{is} a tree; this chapter \emph{builds} the tree when the input is not.  Each
row below is a pointer, not a proof.

\begin{center}
\renewcommand{\arraystretch}{1.25}
\setlength{\tabcolsep}{3pt}
\begin{tabular}{@{}p{0.20\linewidth} p{0.55\linewidth} p{0.21\linewidth}@{}}
\toprule
\textbf{Result} & \textbf{Statement} & \textbf{Where} \\
\midrule
Bridge test & Tree edge \(u\!-\!v\) (child \(v\)) is a bridge iff
\(\mathrm{low}[v]>\tin[u]\). One DFS, \(\bigO{n+m}\). & \lemref{c31-lowchar}, \coderef{c31-bridge} \\
Bridge tree & Delete all bridges, contract 2-edge-components to nodes; bridges are the
tree edges. DP on it. & \defref{c31-twiedge}, \figref{c31-bridgetree} \\
Robbins & A connected graph orients to a strongly connected digraph iff it has no bridge.
CF 118E is one bridge scan. & \thmref{c31-robbins}, \probref{c31-roads} \\
Block pop & Pop one block off the edge stack when \(\mathrm{low}[v]\ge\tin[u]\); the strict
vs.\ non-strict inequality is bridge vs.\ cut vertex. & \thmref{c31-bccpop}, \coderef{c31-bcc} \\
Lone cycle & An edge is on exactly one simple cycle iff its block has \(|E|=|V|\) (rank~1).
CF 962F; brute-force verified. & \lemref{c31-onecycle}, \probref{c31-cycles} \\
Rho structure & A functional graph's component is one directed cycle with in-trees; peel
in-degree \(0\) to isolate the cycles. & \thmref{c31-rhostruct}, \lemref{c31-peel} \\
Cactus & Every block a bridge or a lone cycle \(\Rightarrow\) contracting cycles gives a
tree; DP = tree fold + a ring transfer. & \thmref{c31-cactustree}, \lemref{c31-cycleorder} \\
The pick & Name the redundancy to delete (edge/vertex/arrow/cycle), contract it, reuse
Ch.~12/13; fail all four tests and leave the chapter. & \defref{c31-signature}, \defref{c31-pick} \\
\bottomrule
\end{tabular}
\end{center}

\medskip
{\small\textbf{Algorithms from this chapter.}}
\begin{center}
\renewcommand{\arraystretch}{1.2}
\setlength{\tabcolsep}{2pt}
\begin{tabular}{@{}p{0.33\linewidth} l l >{\raggedright\arraybackslash}p{0.19\linewidth}@{}}
\toprule
\textbf{Algorithm} & \textbf{Prep} & \textbf{Per run} & \textbf{Code} \\
\midrule
Low-link bridges + contraction & \(\bigO{n+m}\) DFS & \(O(1)\) per edge & \coderef{c31-bridge} \\
CF 118E strong orientation & \(\bigO{n+m}\) & one DFS & \coderef{c31-roads} \\
Edge-stack block decomposition & \(\bigO{n+m}\) DFS & \(\bigO m\) total pops & \coderef{c31-bcc} \\
CF 962F lone-cycle edges & \(\bigO{n+m}\) & \(\bigO{m\log m}\) sort & \coderef{c31-cycles} \\
Functional peel + \(\rho,\mu\) & \(\bigO n\) & \(\bigO 1\) per vertex & \coderef{c31-peel} \\
Cactus DP + ring transfer & \(\bigO{n\cdot|S|^2}\) & \(\bigO k\) per cycle & \coderef{c31-cactusdp} \\
CF 835F (Summit) min diameter & \(\bigO n\) cycle + trees & \(\bigO{k\log k}\) windows & \coderef{c31-summit} \\
\bottomrule
\end{tabular}
\end{center}

\chaptersection{Master Decision Flowchart}

Post-study ladder: walk top to bottom and stop at the first branch that matches.  It
decides \emph{whether the graph even has a tree worth extracting} and \emph{which}
redundancy to kill.  Everything failing the top gate leaves the chapter: weights to MST
(\(\chref{c17}\)), enumeration to counting (\(\chref{c32}\)), marked-vertex compression to
virtual trees (\(\chref{c30}\)), directed cuts/flow to graph-land
(\(\secref{c31-boundary}\)'s dominator preview).

\begin{lstlisting}[language={}, basicstyle=\ttfamily\small, numbers=none, frame=single]
  Is the answer a property of ONE given (un)directed graph?
  |
  |-- No: "how many trees/graphs satisfy ..."  -> counting (Ch.30, c32). NOT here.
  |
  |-- Yes, but edges are WEIGHTED and you must MINIMISE total cost -> MST (Ch.31, c17).
  |
  |-- Yes, undirected, question = "does deleting one EDGE disconnect / which edges
  |    are unavoidable / can I orient for round-trips":
  |        +-- BRIDGE TREE (29.1). Contract 2-edge-components; Robbins decides CF 118E.
  |
  |-- Yes, undirected, question = "does deleting one VERTEX disconnect / edges on one
  |    cycle / articulation structure":
  |        +-- BLOCK-CUT TREE (29.2). Edge-stack pops; lone cycle iff |E|=|V| (CF 962F).
  |
  |-- Directed, out-degree EXACTLY ONE (a map f: V->V):
  |        +-- FUNCTIONAL RHO (29.3). Peel in-degree 0, DP on in-trees.
  |
  |-- Undirected, every edge in <= 1 cycle (cycles edge-disjoint):
  |        +-- CACTUS (29.4). Contract cycles to a tree + ring transfer.
  |
  |-- Directed, "every path from s passes through d" (flow cuts):
  |        +-- NOT here. Dominators (preview in 29.5); graph-algorithms territory.
\end{lstlisting}

\chaptersection{Common Mistakes Checklist}

\begin{itemize}[nosep,leftmargin=1.5em]
  \item Skipping the \emph{parent vertex} instead of the \emph{parent edge id} in the
        low-link DFS -- wrong the instant multi-edges exist
        (\(\secref{c31-bridges}\), insight).
  \item Bridge condition \(\mathrm{low}[v]>\tin[u]\) written as \(\ge\) -- \(\ge\) is the
        \emph{block-cut} condition, not bridges; a \(\ge\) bridge flag misses the equality
        case that is a hinge but not an unavoidable edge (\(\lemref{c31-lowchar}\) vs.\
        \(\thmref{c31-bccpop}\)).
  \item Using 2-edge-connected components where blocks are needed (or vice-versa): the two
        decompositions differ exactly on cycles sharing a vertex (CF 962F rung~3).
  \item Forgetting the \(|E|\ge3\) guard, or counting a block's vertices as \(2|E|\)
        instead of distinct endpoints (\(\secref{c31-block}\) warning).
  \item Building a functional graph's children by scanning all of \(f\) -- \(\bigO{n^2}\);
        prebuild reverse adjacency (\(\secref{c31-functional}\) warning~1).
  \item Treating a diamond (two cycles sharing an edge) as a cactus -- the ring transfer is
        invalid there; detect \(|E|>|V|\) at the block pop (\(\secref{c31-cactus}\) warning).
  \item Recursing on a \(3\cdot10^5\) chain without an iterative DFS or a raised stack
        (every section's hardening note).
\end{itemize}

\chaptersection{WA Clinic}

\begin{center}
\renewcommand{\arraystretch}{1.3}
\setlength{\tabcolsep}{3pt}
\begin{tabular}{@{}p{0.30\linewidth} p{0.34\linewidth} p{0.28\linewidth}@{}}
\toprule
\textbf{Symptom} & \textbf{Likely cause} & \textbf{Cite} \\
\midrule
CF 118E prints an orientation but judge says not strong & a bridge slipped through
(\(\ge\) instead of \(>\), or parent-vertex skip) & \(\lemref{c31-lowchar}\) \\
CF 962F counts a whole \(K_4\) block's edges as good & used \(|E|\ge|V|\) not
\(|E|=|V|\) & \(\lemref{c31-onecycle}\) \\
CF 962F wrong on two triangles sharing a vertex & contracted 2-edge-components, not
blocks & \(\secref{c31-block}\) \\
Functional \(\mu\) always \(1\) & kept only the mark-pass, dropped the length rewrite &
\(\coderef{c31-peel}\) \\
Cactus count too high & forgot the cycle's closing edge in \code{ring} &
\(\lemref{c31-cycleorder}\) \\
Summit diameter short by one tree & took \(\min\) over windows but dropped
\(\max_i\) dtree\_i & \(\probref{c31-summit}\) \\
\bottomrule
\end{tabular}
\end{center}

\chaptersection{Quick Reference Card}

\begin{quickref}[One line per tool]
\textbf{Bridge test:} \(u\!-\!v\) tree edge, bridge iff \(\mathrm{low}[v]>\tin[u]\).\\
\textbf{Block pop:} edge stack, pop a block at \(\mathrm{low}[v]\ge\tin[u]\).\\
\textbf{Cut vertex:} non-root \(u\) with a child \(v\), \(\mathrm{low}[v]\ge\tin[u]\); root with
\(\ge2\) tree children.\\
\textbf{Robbins:} strongly orientable \(\iff\) bridgeless.\\
\textbf{Lone-cycle edge:} its block satisfies \(|E|=|V|\) (cycle rank \(1\)).\\
\textbf{Rho:} component \(=\) one cycle + in-trees; peel in-degree \(0\); \(\rho(v)\) tail
length, \(\mu(v)\) cycle length.\\
\textbf{Cactus:} every block bridge-or-cycle; contract cycles \(\to\) tree; DP with a pinned
ring transfer.\\
\textbf{Summit invariant:} break the unique cycle at edge \(j\) \(\Rightarrow\) path diameter
\(\max_{p<q} h_p{+}h_q{+}S_q{-}S_p\); answer \(\max(\max_i\) dtree\(_i,\ \min_j\) window\(_j)\).
\end{quickref}

\chaptersection{Write-From-Memory Drills}

Five skeletons, each \(\le15\) lines, to reproduce without notes.

\begin{lstlisting}[caption={Drill 1 -- low-link bridges and contraction in one pass. Skip the parent edge id, bridge iff low child strictly above tin parent}, label={lst:c31-d1}]
void dfs(int u,int pe){ tin[u]=low[u]=++t;
  for(auto &e:adj[u]){ int v=e.to,id=e.id; if(id==pe) continue;
    if(tin[v]) low[u]=min(low[u],tin[v]);
    else { dfs(v,id); low[u]=min(low[u],low[v]); if(low[v]>tin[u]) bridge[id]=1; } } }
// second DFS: same component unless crossing a bridge
\end{lstlisting}

\begin{lstlisting}[caption={Drill 2 -- CF 118E orientation: tree edge down, back edge up to the ancestor, print 0 on any bridge}, label={lst:c31-d2}]
// inside the DFS: on discovering v (tree) record u->v; on a visited
// ancestor (tin[v]<tin[u]) record u->v; if(low[v]>tin[u]) fail=1;
if(fail) puts("0"); else for each id print stored dir
\end{lstlisting}

\begin{lstlisting}[caption={Drill 3 -- block pop and the lone-cycle test}, label={lst:c31-d3}]
// push edge id on both tree and back(upward) edges
if(low[v]>=tin[u]){ ++stamp; ec=0; distinct=0;
  do { e=pop; ec++; for endpoint a,b: if mark[a]!=stamp mark[a]=stamp, distinct++; } while(e!=id);
  if(ec==distinct && ec>=3) take all ids in block; }
\end{lstlisting}

\begin{lstlisting}[caption={Drill 4 -- functional peel leaving exactly the cycle vertices}, label={lst:c31-d4}]
for v if indeg[v]==0 push; while(q) { v=pop; int u=f[v]; if(--indeg[u]==0) push(u); }
for v if indeg[v]>0 onCycle[v]=1;   // survivors are the union of cycles
\end{lstlisting}

\begin{lstlisting}[caption={Drill 5 -- the Summit's window monoid merge}, label={lst:c31-d5}]
// node covers a consecutive run of cycle positions
mxPlus  = max(H[i]+S[i]);  mxMinus = max(H[i]-S[i]);
best    = max(L.best, R.best, L.mxMinus + R.mxPlus);   // p in L, q in R
// answer window diameter = best over [l, l+k-1]; take min over l, then max with maxDtree
\end{lstlisting}

\chaptersection{Complete Problem Checklist}

\begin{center}
\renewcommand{\arraystretch}{1.2}
\setlength{\tabcolsep}{3pt}
\begin{tabular}{@{}p{0.34\linewidth} >{\raggedright\arraybackslash}p{0.10\linewidth} >{\raggedright\arraybackslash}p{0.22\linewidth} >{\raggedright\arraybackslash}p{0.24\linewidth}@{}}
\toprule
\textbf{Problem} & \textbf{Rating} & \textbf{Tool} & \textbf{Where}\\
\midrule
CF 118E Bertown roads (solved) & 2000 & bridges/Robbins & \(\probref{c31-roads}\)\\
CF 702E Pathes in a Functional Graph (solved) & 2100 & lift on \(f\), (sum,min) & \(\probref{c31-funcpath}\)\\
CF 231E Cactus (solved) & 2100 & contract cycles, LCA fold & \(\probref{c31-cactuscount}\)\\
CF 962F Simple Cycles Edges (solved) & 2400 & blocks/\(|E|{|}{=}|V|\) & \(\probref{c31-cycles}\)\\
CF 835F Roads in the Kingdom (Summit) & 2500 & cycle+window & \(\probref{c31-summit}\)\\
\bottomrule
\end{tabular}
\end{center}

\chaptersection{Recommended Study Order}

Day~1: \(\secref{c31-bridges}\) by hand on the running graph \(G\), then code CF 118E
(\(\probref{c31-roads}\)).  Day~2: \(\secref{c31-block}\), pop the blocks of \(G\), then CF
962F (\(\probref{c31-cycles}\)) -- read the brute-force ledger before trusting
\(\lemref{c31-onecycle}\).  Day~3: \(\secref{c31-functional}\) and \(\secref{c31-cactus}\)
worked traces, then CF 702E (\(\probref{c31-funcpath}\)) and CF 231E
(\(\probref{c31-cactuscount}\)).  Day~4: the Summit \(\probref{c31-summit}\) and the first
four mixed entries; Day~5 the rest of the mixed set and the pick
(\(\defref{c31-pick}\)).

\chaptersection{Difficulty Progression}

\(\lv3\) hand-run the bridge/low test on \(G\); \(\lv4\) code contraction + CF 118E, read
functional rho pieces; \(\lv5\) the edge-stack block pop and CF 962F, the cactus ring
transfer; \(\lv6\) the Summit CF 835F, mixed set at CF~2300--2500.

\chaptersection{Reduction Gallery}

Each row states what the problem \emph{reduces to} after the tree is extracted.

\begin{center}
\renewcommand{\arraystretch}{1.3}
\setlength{\tabcolsep}{3pt}
\begin{tabular}{@{}p{0.28\linewidth} p{0.44\linewidth} p{0.20\linewidth}@{}}
\toprule
\textbf{Problem} & \textbf{Reduction} & \textbf{Result}\\
\midrule
CF 118E & strong orientation \(\to\) bridgeless test \(\to\) emit DFS orientation & \(\bigO{n+m}\) (\(\probref{c31-roads}\))\\
CF 962F & edge on one cycle \(\to\) block is a lone cycle \(\to\) \(|E|=|V|\) & \(\bigO{n+m}\) (\(\probref{c31-cycles}\))\\
CF 231E & paths in a cactus \(\to\) contract cycles, DP on block tree & \(2^{\text{cycles}}\) count (\(\probref{c31-cactuscount}\))\\
CF 702E & iterate a map \(\to\) lift \(f\) with a (sum, min) monoid & \(\bigO{n\log k}\) (\(\probref{c31-funcpath}\))\\
CF 835F & unicyclic, delete one cycle edge \(\to\) min over windows of path diameter & \(\bigO{k\log k}\) (\(\probref{c31-summit}\))\\
\bottomrule
\end{tabular}
\end{center}

\chaptersection{Mixed Problem Set}

Eight entries, each live-verified against the current Codeforces problem set (titles,
ratings and tags re-pulled from the \texttt{problemset.problems} API this session) and
absent from every other roster in the book (checked against \(\sim\!261\) used IDs).
Ordered ascending by rating.  Each hint runs the problem through the graph-to-tree
signature (\(\defref{c31-signature}\)) -- half want a decomposition of this chapter, half
are calibrated rejections that name the true owner.  (Two further entries, CF 702E and
CF 231E, were promoted into the solved labs \(\probref{c31-funcpath}\) and
\(\probref{c31-cactuscount}\).)

\begin{mixedset}
  \practice{Connected Components?}{CF 920E}
    {https://codeforces.com/problemset/problem/920/E}{CF 2100 L5}
    {Tags \emph{data structures, dfs and similar, dsu, graphs}. Find the connected components
     of the \emph{complement} of a given sparse graph.  A calibration rejection: no edge of
     the original graph is special in the connectivity sense -- there is no bridge or
     articulation to extract, and the residual object is not a tree.  Owner: DSU / BFS on
     the complement with the standard ``remaining set'' trick (\(\chref{c28}\)'s family).
     \(\defref{c31-signature}\) fails T1 (deleting a given edge does not decompose it), so
     do not reach for the bridge tree; the graph is dense by negation.}
  \practice{Network Safety}{CF 1039C}
    {https://codeforces.com/problemset/problem/1039/C}{CF 2200 L5}
    {Tags \emph{dfs and similar, dsu, graphs, math, sortings}. Add the cheapest extra cable so
     that after \emph{any single road} fails the network is still connected -- literally
     asking which roads are bridges (\(\secref{c31-bridges}\)).  T1 fires: the answer turns
     on edge-robustness.  Find the bridges with the low-link DFS (\(\lemref{c31-lowchar}\));
     a non-bridge road already has an alternative route, so only the bridge structure
     matters, and the minimum non-bridge edge weight is the candidate.  One DFS plus a scan.}
  \practice{Interestring graph and Apples}{CF 9E}
    {https://codeforces.com/problemset/problem/9/E}{CF 2300 L5}
    {Tags \emph{dfs and similar, dsu, graphs}. A game where a set of edges is kept such that
     the chosen edges stay connected while edges are removed; the winner depends on how many
     \emph{independent} (cycle-breaking) edges remain.  The crux is the DFS-tree edge
     classification of \(\secref{c31-bridges}\): an edge you may delete without splitting
     the graph is precisely one on a cycle (not a bridge), and the count of such ``free''
     edges is \(m-(n-\text{components})\).  Signature~T1: it is bridge-avoidance in game
     clothing -- run the low-link once and reason about the bridge count, do not simulate
     the game.}
  \practice{Tourist Reform}{CF 732F}
    {https://codeforces.com/problemset/problem/732/F}{CF 2300 L5}
    {Tags \emph{dfs and similar, graphs}. Orient the roads to maximise the number of edges
     you can traverse both ways along some closed route -- the other face of CF 118E.  Inside
     a 2-edge-connected component you can orient every edge to lie on a directed cycle
     (Robbins, \(\thmref{c31-robbins}\)); a bridge is forced one-way and is the only place
     you lose.  Contract the bridge tree, orient each component strongly, and bridges
     contribute nothing.  Signature~T1 + the bridge contraction of \(\secref{c31-bridges}\).}
  \practice{Graph Traveler}{CF 1200F}
    {https://codeforces.com/problemset/problem/1200/F}{CF 2300 L5}
    {Tags \emph{brute force, data structures, dfs and similar, dp, graphs, implementation,
     math, number theory}. Despite the ``graph'' title the constraint is an ordering/DP
     problem, not a decomposition: signature~T4 fails because contracting redundancy does
     not leave a tree you DP on -- the object is a travel sequence.  Route it to tree/graph
     DP over a fixed traversal (\(\chref{c10}\)) with the number-theoretic bookkeeping the
     tags flag; use this entry to practise \emph{saying no} to the bridge reflex.}
  \practice{Graph Cutting}{CF 2231E}
    {https://codeforces.com/problemset/problem/2231/E}{CF 2300 L5}
    {Tags \emph{dfs and similar, dp, math, trees}. Decide feasibility of cutting the graph by
     articulation-style decomposition; the answer is governed by which vertices are hinges,
     so T1 (vertex-deletion) points at the block-cut tree (\(\secref{c31-block}\)): pop
     blocks with the edge stack and DP over the block-cut structure.  The trap is treating
     it as an edge-bridge problem -- it is the \emph{vertex} decomposition you want
     (\(\thmref{c31-bccpop}\)), and a block that is a lone cycle behaves differently from a
     hinge.}
  \practice{Masha and Cactus}{CF 856D}
    {https://codeforces.com/problemset/problem/856/D}{CF 2400 L5}
    {Tags \emph{dp, trees}. Path cost queries over a weighted cactus with edge additions --
     the cactus of \(\secref{c31-cactus}\) at full strength.  Every block is a bridge or a
     cycle, so contract cycles, and the tree of blocks supports the DP / heavy-path answers
     exactly as a tree would (\(\thmref{c31-cactustree}\)), with a ring transfer at each
     cycle node.  Signature~T3 + Ch.~13 on the contracted tree; the whole point is that a
     cactus is ``a tree of small cycles'', so tree machinery survives.}
  \practice{AND Graph}{CF 986C}
    {https://codeforces.com/problemset/problem/986/C}{CF 2500 L5}
    {Tags \emph{bitmasks, dfs and similar, dsu, graphs}. Vertices are integers; two are
     adjacent by a rule on their bitwise AND -- a dense graph given implicitly.  The
     maximal-components question has no bridge or articulation to extract (every bit
     position creates many overlapping cycles), so signature~T1 fails and there is no
     tree-of-components; the correct move is DFS/DSU over the \emph{bitmask} structure
     (\(\chref{c28}\)).  Take this as the set's cleanest rejection at top rating: the graph
     is connected in ways no single edge deletion changes.}
\end{mixedset}

% =====================================================================
%  STATION 12: THE SUMMIT -- CF 835F Roads in the Kingdom
% =====================================================================
\chaptersection{The Summit -- CF 835F Roads in the Kingdom}

\begin{notebox}[Why this is the chapter top rope]
Everything before it decomposed a graph into a tree of pieces.  The Summit is
the dual: the graph is already \emph{almost} a tree -- \(n\) vertices, \(n\) edges,
one cycle -- and you must delete exactly one edge so what remains is a genuine
tree, with the goal of minimising that tree's diameter (the statement's
``inconvenience'').  Three chapter tools do the heavy lifting: the leaf-peeling
cycle isolation of \(\secref{c31-functional}\) (\(\lemref{c31-peel}\)), the
``contract the cycle and DP on the tree of pieces'' mindset of
\(\secref{c31-cactus}\), and the low-link/bridge view of \(\secref{c31-bridges}\)
that decides which edges are legal to cut.  Two earlier chapters close the gap:
Ch.~12 (\(\chref{c09}\)) gives the two-sweep height/diameter fold that runs on
each attached tree, and Ch.~19 (\(\chref{c13}\)) supplies the segment-tree monoid
that turns the cycle-window scan into \(\bigO{k\log k}\).  Both official samples
are replayed by hand below: \(n{=}3\to 5\) and \(n{=}5\to 18\).
\end{notebox}

\solvedproblem{c31-summit}{Roads in the Kingdom}{CF 835F}
  {https://codeforces.com/problemset/problem/835/F}{CF 2500 L5}
  {peel leaves to isolate the unique cycle (\(\lemref{c31-peel}\)); for each cycle
   vertex \(c_i\) compute the attached tree's height \(h_i\) and internal diameter
   \(\mathit{dtree}_i\) with the two-largest-depth fold (Ch.~12, \(\chref{c09}\));
   walk the cycle in order \(c_0,\ldots,c_{k-1}\) with edge lengths
   \(w_0,\ldots,w_{k-1}\); cutting cycle edge \(w_j\) turns the cycle into a path on
   positions \([j{+}1,\,j{+}k]\); the path's through-diameter is
   \(\max_{a<b}((h_a-S_a)+(h_b+S_b))\) on that window; a doubled prefix and a
   segment tree with the monoid triple (mxMinus, mxPlus, best) answers each window
   in \(\bigO{\log k}\); the final answer is
   \(\max(\max_i\mathit{dtree}_i,\ \min_j\mathrm{window}_j)\)}

\begin{problemstatement}[Roads in the Kingdom -- statement, constraints and both samples verified against the live CF page]
The Kingdom has \(n\) towns and \(n\) bidirectional roads.  Road \(i\) joins towns
\(u_i,v_i\) and has length \(\ell_i\).  No road connects a town to itself, no two roads
join the same pair, and the graph is connected -- so \(G\) is unicyclic (a tree with
exactly one extra edge, forming a unique simple cycle).

Call the \textbf{inconvenience} of the road network the maximum, over all pairs of
towns, of their shortest-path distance.  The king closes exactly one road, subject
to the condition that the remaining network is still connected; you must choose
that road to \emph{minimise} the inconvenience.

\textbf{Input.} \(n\) on the first line (\(3\le n\le 2\cdot 10^{5}\)); then \(n\) lines,
each \(u_i\ v_i\ \ell_i\) with \(1\le u_i,v_i\le n\) and \(1\le \ell_i\le 10^{9}\).
\textbf{Output.} A single integer -- the minimum possible inconvenience.  Time limit
2\,s, memory 256\,MB.

\textbf{Samples (both official).}  \#1: \(n{=}3\) with roads
\((1,2,4), (2,3,5), (1,3,1)\) \(\to\) \texttt{5}.  \#2: \(n{=}5\) with roads
\((2,3,7), (3,1,9), (4,1,8), (3,5,4), (4,5,5)\) \(\to\) \texttt{18}.
\end{problemstatement}

\begin{figure}[H]
\centering
\begin{tikzpicture}[diagram, scale=1.0, transform shape,
   vtx/.style={circle, draw=primaryblue!70, fill=primaryblue!10, line width=0.8pt,
               minimum size=7mm, font=\scriptsize\bfseries, inner sep=1pt}]
  \node[vtx] (v1) at (0,1.4) {1};
  \node[vtx] (v2) at (-1.4,-0.4) {2};
  \node[vtx] (v3) at (1.4,-0.4) {3};
  \draw[warningred, line width=1.4pt] (v1)--(v2) node[midway,above left,font=\scriptsize\bfseries] {4};
  \draw[primaryblue!70, line width=1.1pt] (v2)--(v3) node[midway,below,font=\scriptsize\bfseries] {5};
  \draw[primaryblue!70, line width=1.1pt] (v3)--(v1) node[midway,above right,font=\scriptsize\bfseries] {1};
  \node[annot, font=\scriptsize, anchor=north, text width=10cm] at (0,-1.15)
    {cut 1-2 gives broken path 2-3-1 with length 5+1=6;  cut 2-3 gives 4+1=5;  cut 3-1 gives 4+5=9};
\end{tikzpicture}
\caption{The Summit's official sample \(n{=}3\): a single triangle with no attached
trees, so every \(h_i{=}0\) and every \(\mathit{dtree}_i{=}0\).  The answer is the pure
cycle-window minimum \(\min(6,5,9){=}5\).}
\label{fig:c31-summit}
\end{figure}

\begin{attemptgate}[Attempt it yourself first -- \lv{5}, budget a real sitting]
Three questions before any code: (i) which edges are legal to close -- what does
the statement's ``still connected'' condition force on the edge you pick?
(ii) after cutting such an edge, is the result a graph or a tree, and what is its
diameter a function of?  (iii) the cycle has length \(k\); how does the tree's
diameter depend on \emph{which} cycle edge you cut?
\peek{The graph has \(n\) vertices, \(n\) edges, and is connected -- exactly one
cycle.  An edge whose removal preserves connectivity is \emph{not} a bridge
(\(\lemref{c31-lowchar}\)), so the legal cuts are exactly the \(k\) cycle edges.  After
cutting one you have a genuine spanning tree, and the inconvenience is that tree's
ordinary diameter -- Ch.~12 territory.}
\nudge{The spanning tree obtained by cutting cycle edge \(w_j\) has diameter
\(\max(\max_i\mathit{dtree}_i,\ \max_{a<b}h_a+\mathrm{dist}_j(a,b)+h_b)\).  The first
argument does not depend on the cut, so precompute it once.  The second term
looks \emph{linear} in the broken path: writing \(S\) for the running prefix of the
cycle edge lengths, \(\mathrm{dist}_j(a,b)=S_b-S_a\) for positions \(a<b\) along the
broken path, so the through-cycle term becomes
\(\max_{a<b}((h_a-S_a)+(h_b+S_b))\) on a length-\(k\) window.}
\fullhint{Lay the cycle out twice as \(c_0,\ldots,c_{k-1},c_0,\ldots,c_{k-1}\) and
build a doubled prefix \(S_0{=}0\), \(S_t{=}S_{t-1}{+}w_{(t-1)\bmod k}\).  Cutting
edge \(w_j\) leaves the broken path in order \(c_{j+1},\ldots,c_{j+k}\), so the
length-\(k\) window is \([L, L{+}k{-}1]\) with \(L{=}j{+}1\); the answer to the
through-cycle term is \(\max_{a<b}((h_a-S_a)+(h_b+S_b))\) over that window.  Build
a segment tree on the doubled array whose node over a consecutive run stores a
triple \((\mathrm{mxMinus}{=}\max(h{-}S),\ \mathrm{mxPlus}{=}\max(h{+}S),\ \mathrm{best}{=}\max_{a<b}(h_a{-}S_a{+}h_b{+}S_b))\);
merge two children by \(\mathrm{best}{=}\max(\mathrm{L.best},\mathrm{R.best},\mathrm{L.mxMinus}{+}\mathrm{R.mxPlus})\).
Query each window in \(\bigO{\log k}\); final answer
\(=\max(\max_i\mathit{dtree}_i,\ \min_j\mathrm{window}_j)\).}
\end{attemptgate}

\solstep{Understanding}
The graph is unicyclic: \(n\) vertices, \(n\) edges, connected.  The bridge/block
machinery of 29.1--29.2 tells you exactly which edges are on the cycle -- those
and only those are legal to close; every off-cycle edge is a bridge and deleting
it disconnects the kingdom.  Two facts collapse the problem.  (1) The unique
cycle is isolated in \(\bigO n\) by peeling leaves -- the same trick as
\(\lemref{c31-peel}\), no low-link needed.  (2) Each cycle vertex \(c_i\) is the root
of an ``attached tree'' hanging off the cycle; a Ch.~12 fold reports that tree's
height \(h_i\) (deepest distance from \(c_i\)) and internal diameter
\(\mathit{dtree}_i\).  Cutting cycle edge \(w_j\) turns the cycle into a simple path;
every tree path is either inside an attached tree (bounded by
\(\max_i\mathit{dtree}_i\)) or it leaves \(c_a\), walks along the broken path to
\(c_b\), and climbs into \(c_b\)'s tree -- length at most
\(h_a+\mathrm{dist}_j(a,b)+h_b\).  The task reduces to a path-diameter formula
queried \(k\) times on a doubled array, and that is exactly where the
segment-tree monoid of Ch.~19 earns its keep.

\solstep{Key Observations}
\begin{observation}
\textbf{Legal cuts are the cycle edges; everything else disconnects.}
The graph has \(n\) vertices, \(n\) edges, is connected, and admits at least one
cycle-edge removal that preserves connectivity (the problem's own guarantee).
An edge whose removal preserves connectivity is not a bridge
(\(\lemref{c31-lowchar}\)); in a unicyclic graph the non-bridge edges are precisely
those on the unique cycle.  So the feasible set has size \(k\) and leaf peeling
(\(\lemref{c31-peel}\)) finds it in \(\bigO n\) -- no DFS, no low-link pass needed.
\end{observation}
\begin{observation}
\textbf{The diameter splits into a fixed part and a cut-dependent part.}
Fix a cut \(j\).  Every simple path in the resulting spanning tree is either
(i)~inside one attached tree \(A_i\), contributing at most \(\mathit{dtree}_i\); or
(ii)~leaves \(A_a\) at \(c_a\), walks along the broken-cycle path to \(c_b\), and
enters \(A_b\).  Its length is at most \(h_a+\mathrm{dist}_j(a,b)+h_b\) and is
achieved at the deepest vertices of \(A_a\) and \(A_b\).  Therefore
\(\mathrm{diam}(T_j)=\max(\max_i\mathit{dtree}_i,\ \max_{a<b}(h_a+\mathrm{dist}_j(a,b)+h_b))\),
and because the first term does not depend on \(j\),
\[
\text{answer}=\max\!\Bigl(\max_i\mathit{dtree}_i,\ \min_j\max_{a<b\in\mathrm{win}(j)}(h_a+\mathrm{dist}_j(a,b)+h_b)\Bigr).
\]
Writing \(\mathrm{dist}_j(a,b)=S_b-S_a\) turns the inner sum into
\((h_a-S_a)+(h_b+S_b)\), a two-sided range maximum the segment tree handles directly.
\end{observation}
\begin{observation}
\textbf{The window is length \(k\); the doubled array is length \(2k\).}
Cutting edge \(w_j\) leaves the cycle path \(c_{j+1},c_{j+2},\ldots,c_{j+k}\) (indices
mod \(k\)).  Doubling \((c_0,\ldots,c_{k-1},c_0,\ldots,c_{k-1})\) makes this window a
contiguous range \([L,\,L{+}k{-}1]\) with \(L=j+1\in[1,k]\), all inside \([0,2k{-}1]\).
The recurrence \(S_t=S_{t-1}+w_{(t-1)\bmod k}\) gives the running path distance
from position~0, and \(\mathrm{dist}_j(a,b)=S_b-S_a\) is measured along the \emph{broken}
path -- not the shorter arc.  Skipping the doubling and using modular indices
\emph{inside} the formula is dead end~3 below.
\end{observation}

\solstep{Think It With Me}
\begin{thinkbox}
\deadend{Try every edge, remove it, BFS/DFS the remaining graph, take its diameter}
{Two failures at once.  \emph{Correctness}: an edge not on the cycle is a bridge,
and removing it disconnects the graph -- a naive BFS reading max distance from
the start component sees only that component and reports a number that ignores
the disconnected towns entirely.  \emph{Performance}: \(n\) cuts \(\times\) \(\bigO n\)
BFS is \(\bigO{n^2}=4\cdot 10^{10}\) at \(n=2\cdot 10^{5}\), which is \(\tle\) even
before considering the wrongness.  In sample 2 the road \((2,3,7)\) is a bridge
(town 2 is a leaf off the cycle); removing it leaves town 2 unreachable, but a
BFS from town 1 reports a finite number that silently pretends town 2 vanished.}
\deadend{Cut the longest cycle edge greedily -- the longest edge removed should
shrink the longest path}
{Plausible, wrong.  In sample 2 the unique cycle runs through towns 1, 3, 5, 4
with edges \((1\!-\!3,\,9)\), \((3\!-\!5,\,4)\), \((5\!-\!4,\,5)\), \((4\!-\!1,\,8)\).
The longest cycle edge is the 9 joining 1--3; cutting it gives a tree whose
diameter is 24 (see the sample-2 tracebox below).  The official answer 18 is
achieved by cutting the edge \(4\!-\!1\) of length 8, not the longest.  Greedy
overshoots by 6 on the very sample that ships with the problem.  \textbf{WA}.}
\deadend{Compute only the through-cycle term and skip \(\max_i\mathit{dtree}_i\);
attached trees contribute nothing beyond their heights}
{A diameter path can lie entirely \emph{inside} one attached tree, and the
through-cycle formula never sees it.  Construct a small case: a triangle
\(c_0c_1c_2\) with edges all of length 1, and let \(c_0\) hang a two-legged tree
\(c_0\!-\!u\) of length 1, then \(u\!-\!b\) and \(u\!-\!d\) of length 50 each.  Then
\(h_0{=}51\) (the deepest vertex from \(c_0\)), but \(\mathit{dtree}_0{=}100\) (the
path \(b\!-\!u\!-\!d\)).  Every window's through-cycle term is at most
\(51+2+0{=}53\); the height-only formula prints 53, the true tree diameter on
every legal cut is 100 (the attached path dominates).  \textbf{WA}, and a very
common one -- Ch.~12's diameter formula says a tree's diameter is either
through the root or entirely inside one child's subtree, and only the second
is missed by the through-cycle term.}
\aha{Split the work along the two-term decomposition.  \textbf{Phase~1} (linear):
peel leaves to isolate the cycle in \(\bigO n\) (\(\lemref{c31-peel}\)); walk the
cycle in order to fix arrays \(c_0..c_{k-1}\) and \(w_0..w_{k-1}\); per cycle vertex
run a bounded DFS through non-cycle vertices only and compute \((h_i,\mathit{dtree}_i)\)
with the two-largest-depth fold; hold \(\mathit{dtreeBig}{=}\max_i\mathit{dtree}_i\).
\textbf{Phase~2} (near-linear in \(k\)): lay the cycle out twice, build the doubled
prefix \(S\), form per-position values \(\mathit{valMinus}_t{=}h_{t\bmod k}{-}S_t\) and
\(\mathit{valPlus}_t{=}h_{t\bmod k}{+}S_t\); feed a segment tree whose node stores the
triple \((\mathrm{mxMinus},\mathrm{mxPlus},\mathrm{best})\) and merge by
\(\mathrm{best}{=}\max(\mathrm{L.best},\mathrm{R.best},\mathrm{L.mxMinus}{+}\mathrm{R.mxPlus})\); query
each window \([L,L{+}k{-}1]\) for \(L{=}1..k\) in \(\bigO{\log k}\).  Answer
\(=\max(\mathit{dtreeBig},\min_L\mathrm{query}_L.\mathrm{best})\).}
\end{thinkbox}

\solstep{The Approach Ladder}
\begin{center}\renewcommand{\arraystretch}{1.25}\setlength{\tabcolsep}{3pt}
\begin{tabular}{@{}c p{0.24\linewidth} p{0.28\linewidth} l c@{}}
\toprule
\textbf{\#} & \textbf{Approach} & \textbf{Idea} & \textbf{Time} & \textbf{Verdict}\\
\midrule
1 & \lv{2}\ try every edge, BFS each result & ignores which cuts are legal; disconnected silently under-reports & \(\bigO{n^2}\) & \tle \\
2 & \lv{3}\ every cycle edge, BFS each diameter & correct candidate set, no structure used & \(\bigO{kn}\) & \tle \\
3 & \lv{4}\ greedy: cut the longest cycle edge & plausible heuristic, fails on sample 2 by 6 & \(\bigO n\) & \textcolor{warningred}{\ding{55}\ \textbf{WA}} \\
4 & \lv{5}\ peel, cycle walk, attached h and dtree, doubled prefix, segtree monoid & two-term split + range-max over pairs & \(\bigO{n+k\log k}\) & \ac \\
\bottomrule
\end{tabular}\end{center}
Rungs 1--2 ignore the chapter's structure; rung~3 is the greedy temptation that
Robbins and the bridge tree exist to punish; rung~4 is the chapter's thesis
assembled: identify the redundancy (the unique cycle), cut it exactly once, and
let a segment-tree monoid enumerate all \(k\) possibilities at \(\bigO{\log k}\) each.

\solstep{Solution}
Peel leaves: push every degree-1 vertex on a queue, mark each popped vertex off
the cycle, decrement its still-on-cycle neighbours -- survivors are exactly the
cycle vertices in \(\bigO n\) (\(\lemref{c31-peel}\)).  Walk the cycle in order: pick
the first on-cycle vertex \(c_0\), repeatedly step to the unique cycle neighbour
different from the predecessor, record each edge length, stop when you return to
\(c_0\); the result is \(c_0..c_{k-1}\) and \(w_0..w_{k-1}\).  For each \(c_i\) run a
bounded DFS that visits only non-cycle vertices and returns
\((h_i,\mathit{dtree}_i)\) via the two-largest-depth fold; track
\(\mathit{dtreeBig}=\max_i\mathit{dtree}_i\).  Build a doubled prefix
\(S_0{=}0, S_t{=}S_{t-1}{+}w_{(t-1)\bmod k}\) for \(t\in[0,2k{-}1]\); per position
\(t\) store \(\mathit{valMinus}_t{=}h_{t\bmod k}{-}S_t\) and
\(\mathit{valPlus}_t{=}h_{t\bmod k}{+}S_t\).  Build a segment tree whose node is
\((\mathrm{mxMinus},\mathrm{mxPlus},\mathrm{best})\), leaves carry \(\mathrm{best}{=}-\infty\),
and the merge takes the max of the three cross possibilities.  Query each window
\([L,L{+}k{-}1]\) for \(L{=}1..k\); the answer is
\(\max(\mathit{dtreeBig},\min_L\mathrm{query}_L.\mathrm{best})\).

\solstep{Proof of Correctness}
\begin{proofbox}
\textbf{(Legal cuts are exactly the cycle edges.)}  In a connected graph with
\(n\) vertices and \(n\) edges, the cycle rank is \(|E|-|V|+{1}=1\), so the graph is
unicyclic: one simple cycle \(C\), and every edge not on \(C\) is a bridge
(\(\lemref{c31-lowchar}\)).  Deleting a bridge disconnects; deleting a cycle edge
breaks the cycle and leaves a spanning tree.  So the feasible set equals
\(C\)'s edges, and leaf peeling finds it (\(\lemref{c31-peel}\)).

\textbf{(Diameter decomposition.)}  Fix a cut \(j\); the resulting spanning tree
is \(T_j\).  Every simple path \(P\) in \(T_j\) either (a)~lies inside one attached
tree \(A_i\), contributing at most \(\mathit{dtree}_i\); or (b)~starts at a deepest
vertex of \(A_a\), passes through \(c_a\), follows the broken-cycle path to \(c_b\),
and climbs to a deepest vertex of \(A_b\), contributing
\(h_a+\mathrm{dist}_{P_j}(a,b)+h_b\).  These two cases cover every simple path
(a path through one cycle vertex's tree and back to the same tree falls in
(a); any path visiting two cycle vertices falls in (b)), so
\(\mathrm{diam}(T_j)=\max(\max_i\mathit{dtree}_i,\ \max_{a<b}(h_a+\mathrm{dist}_j(a,b)+h_b))\).

\textbf{(Path form via doubled prefix.)}  Cutting edge \(w_j\) leaves the path
\(c_{j+1},\ldots,c_{j+k}\).  Doubling and setting \(S_t=\sum_{a<t}w_{a\bmod k}\) gives
\(\mathrm{dist}_{P_j}(a,b)=S_b-S_a\) for the window positions
\(j+1\le a<b\le j+k\); therefore the through-cycle term equals
\(\max_{a<b}((h_a-S_a)+(h_b+S_b))\) over that window.

\textbf{(Monoid correctness.)}  A node covering a run \(R\) stores
\(\mathrm{mxMinus}=\max_{t\in R}\mathit{valMinus}_t\),
\(\mathrm{mxPlus}=\max_{t\in R}\mathit{valPlus}_t\),
\(\mathrm{best}=\max_{a<b,\,a,b\in R}(\mathit{valMinus}_a+\mathit{valPlus}_b)\).
For a merge \(R=L\cup\mathrm{Rt}\) with every \(L\)-index less than every
\(\mathrm{Rt}\)-index, pairs in \(R\) split three ways: both in \(L\) (covered by
\(L.\mathrm{best}\)), both in \(\mathrm{Rt}\) (covered by \(\mathrm{Rt}.\mathrm{best}\)),
or \(a\in L, b\in\mathrm{Rt}\) whose maximum value is
\(L.\mathrm{mxMinus}+\mathrm{Rt}.\mathrm{mxPlus}\).  Leaves carry
\(\mathrm{best}{=}-\infty\) (a singleton has no pair).  Standard segment-tree
induction gives that a range query composes exactly the three quantities over
the queried range, so \(\mathrm{query}(L,R).\mathrm{best}\) equals the true
window maximum.

\textbf{(Answer.)}  Minimising over \(j\) and combining with the \(j\)-independent
\(\max_i\mathit{dtree}_i\) via \(\max\) yields the required minimum inconvenience,
since \(j\) ranges over exactly the feasible cuts and
\(\max(\mathrm{const},\min_j\mathrm{win}_j)=\min_j\max(\mathrm{const},\mathrm{win}_j)\)
when the first argument does not depend on \(j\).  \(\Box\)
\end{proofbox}

\solstep{Implementation}
\begin{lstlisting}[caption={CF 835F Summit: peel leaves to isolate the cycle, walk the cycle in order, compute per-vertex attached-tree height and diameter, then a doubled prefix plus a segment tree over triples (mxMinus, mxPlus, best) answers the min-diameter window query. ASCII identifiers, no underscores.}, label={lst:c31-summit}]
#include <bits/stdc++.h>
using namespace std;
typedef long long ll;
const ll INF = (1LL<<62);
struct Edge{ int to; ll len; };
int main(){
  int n; if(scanf("%d",&n)!=1) return 0;
  vector<vector<Edge> > adj(n);
  for(int i=0;i<n;i++){ int u,v; ll w; scanf("%d%d%lld",&u,&v,&w); --u; --v;
    adj[u].push_back((Edge){v,w}); adj[v].push_back((Edge){u,w}); }
  // Phase 1a: peel leaves; survivors are exactly the cycle vertices
  vector<int> degCnt(n); vector<char> onCycle(n,1);
  queue<int> Q;
  for(int i=0;i<n;i++){ degCnt[i]=(int)adj[i].size(); if(degCnt[i]==1) Q.push(i); }
  while(!Q.empty()){ int u=Q.front(); Q.pop(); onCycle[u]=0;
    for(size_t j=0;j<adj[u].size();j++){ int v=adj[u][j].to;
      if(onCycle[v]){ if(--degCnt[v]==1) Q.push(v); } } }
  // Phase 1b: walk the cycle in order, record edge lengths
  int start=0; while(!onCycle[start]) start++;
  vector<int> cyc; vector<ll> edgeLen;
  int prev=-1, cur=start;
  while(true){ cyc.push_back(cur); int nxt=-1; ll wl=0;
    for(size_t j=0;j<adj[cur].size();j++){ int v=adj[cur][j].to;
      if(onCycle[v] && v!=prev){ nxt=v; wl=adj[cur][j].len; break; } }
    edgeLen.push_back(wl); prev=cur; cur=nxt;
    if(cur==start) break; }
  int k=(int)cyc.size();
  // Phase 1c: per cycle vertex, attached-tree height and diameter via two-max fold
  function<pair<ll,ll>(int,int)> dfsAttach = [&](int u,int p) -> pair<ll,ll> {
    ll b1=0, b2=0, diam=0;
    for(size_t j=0;j<adj[u].size();j++){ int v=adj[u][j].to; ll w=adj[u][j].len;
      if(v!=p && !onCycle[v]){ pair<ll,ll> r=dfsAttach(v,u);
        if(r.second>diam) diam=r.second;
        ll d=r.first+w;
        if(d>=b1){ b2=b1; b1=d; } else if(d>b2){ b2=d; } } }
    if(b1+b2>diam) diam=b1+b2;
    return make_pair(b1, diam);
  };
  vector<ll> H(k,0); ll dtreeBig=0;
  for(int i=0;i<k;i++){ pair<ll,ll> r=dfsAttach(cyc[i],-1);
    H[i]=r.first; if(r.second>dtreeBig) dtreeBig=r.second; }
  // Phase 2a: doubled prefix and per-position values
  vector<ll> Pref(2*k,0);
  for(int t=1;t<2*k;t++) Pref[t]=Pref[t-1]+edgeLen[(t-1)%k];
  vector<ll> valM(2*k), valP(2*k);
  for(int t=0;t<2*k;t++){ valM[t]=H[t%k]-Pref[t]; valP[t]=H[t%k]+Pref[t]; }
  // Phase 2b: segment tree monoid triple (mxMinus, mxPlus, best)
  struct Node{ ll mxMinus, mxPlus, best; };
  function<Node(Node,Node)> mergeNode = [](Node a, Node b){
    Node r;
    r.mxMinus = a.mxMinus>b.mxMinus?a.mxMinus:b.mxMinus;
    r.mxPlus  = a.mxPlus>b.mxPlus?a.mxPlus:b.mxPlus;
    ll cross  = a.mxMinus + b.mxPlus;
    ll inner  = a.best>b.best?a.best:b.best;
    r.best    = inner>cross?inner:cross;
    return r;
  };
  vector<Node> seg(8*k+8);
  function<void(int,int,int)> build = [&](int p,int l,int r){
    if(l==r){ seg[p]=(Node){valM[l], valP[l], -INF}; return; }
    int m=(l+r)>>1; build(p<<1,l,m); build(p<<1|1,m+1,r);
    seg[p]=mergeNode(seg[p<<1], seg[p<<1|1]); };
  function<Node(int,int,int,int,int)> query = [&](int p,int l,int r,int ql,int qr) -> Node{
    if(ql<=l && r<=qr) return seg[p];
    int m=(l+r)>>1;
    if(qr<=m) return query(p<<1,l,m,ql,qr);
    if(ql>m)  return query(p<<1|1,m+1,r,ql,qr);
    return mergeNode(query(p<<1,l,m,ql,qr), query(p<<1|1,m+1,r,ql,qr)); };
  build(1,0,2*k-1);
  ll winMin=INF;
  for(int L=1;L<=k;L++){ Node res=query(1,0,2*k-1,L,L+k-1);
    if(res.best<winMin) winMin=res.best; }
  printf("%lld\n", dtreeBig>winMin?dtreeBig:winMin);
  return 0;
}
\end{lstlisting}

\solstep{Trace on the Official Samples}
\begin{tracebox}[Sample 1: \(n{=}3\) triangle -- every \(h_i{=}0\), a pure cycle-window minimum]
All three towns have degree~2, so nothing is peeled and every vertex is on the
cycle.  Walking from town~1 (0-based 0) gives \(\mathit{cyc}{=}[0,1,2]\) and
\(\mathit{edgeLen}{=}[4,5,1]\) with \(k{=}3\).  Each \(\mathit{dfsAttach}\) meets only cycle
neighbours, so \(H{=}[0,0,0]\) and \(\mathit{dtreeBig}{=}0\).
Doubled prefix \(S=[0,4,9,10,14,19]\); \(\mathit{valMinus}=-S\), \(\mathit{valPlus}=S\).
Windows of length \(k{=}3\) at \(L{=}1,2,3\), best \(=\max_{a<b}(S_b-S_a)=S_{\text{last}}-S_{\text{first}}\):
\(L{=}1\) positions \(\{1,2,3\}\): \(S_3{-}S_1=10{-}4=6\).
\(L{=}2\) positions \(\{2,3,4\}\): \(S_4{-}S_2=14{-}9=5\).
\(L{=}3\) positions \(\{3,4,5\}\): \(S_5{-}S_3=19{-}10=9\).
\(\min_j\mathrm{win}_j=\min(6,5,9)=5\).  Answer \(=\max(0,5)=\boxed{5}\) -- matches the
official output.  The three windows are exactly the three roads you can close
\(\figref{c31-summit}\). \(\Box\)
\end{tracebox}

\begin{tracebox}[Sample 2: \(n{=}5\) -- peel isolates a 4-cycle, one attached leaf at town 3]
0-based edges \((1,2,7),(2,0,9),(3,0,8),(2,4,4),(3,4,5)\).  Degrees \((2,1,3,2,2)\);
only town~1 (1-based town~2) has degree~1, so it is the sole peel seed: popping it
sets \(\mathit{onCycle}[1]{=}0\) and drops town~2's degree to~2 -- no new seed.  Cycle
vertices are towns \(0,2,3,4\) (1-based \(1,3,4,5\)).  Walking from 0 gives
\(\mathit{cyc}{=}[0,2,4,3]\), \(\mathit{edgeLen}{=}[9,4,5,8]\), \(k{=}4\).  The bounded DFS at
\(\mathit{cyc}[1]{=}2\) reaches the off-cycle town~1 on a length-7 edge, so \(H[1]{=}7\),
\(\mathit{dtree}_1{=}7\); the other cycle vertices have no off-cycle neighbours.
Thus \(H{=}[0,7,0,0]\), \(\mathit{dtreeBig}{=}7\).
\(S=[0,9,13,18,26,35,39,44]\), so \(\mathit{valMinus}=[0,-2,-13,-18,-26,-28,-39,-44]\) and
\(\mathit{valPlus}=[0,16,13,18,26,42,39,44]\) (only index~1 and~5 carry the \(+7\)).
Length-4 windows \(L{=}1..4\): \(\mathrm{win}_1{=}24,\ \mathrm{win}_2{=}29,\ \mathrm{win}_3{=}24,\ \mathrm{win}_4{=}18\)
(each the max of \(\mathit{valMinus}_a{+}\mathit{valPlus}_b\) over \(a<b\) in the window; the
value 18 at \(L{=}4\) is the pair \((4,7)\): \(-26{+}44=18\)).
\(\min_j\mathrm{win}_j=18\).  Answer \(=\max(7,18)=\boxed{18}\) -- matches.  Greedy
rung~3 (close the longest cycle road, length 9) is \(L{=}1\to 24\), off by~6: dead
end~2 reproduced on the official data.  The attached town~1 (\(\mathit{dtree}\)~7) never
binds here because the cycle path already exceeds it -- but a bushy attached tree
would make \(\mathit{dtreeBig}\) the answer, which is why the \(\max\) stays. \(\Box\)
\end{tracebox}

\solstep{Complexity}
\begin{complexitybox}
\textbf{Time} \(\bigO{n + k\log k}\).  Peeling is \(\bigO n\); the cycle walk touches
each cycle edge once, \(\bigO k\); across all \(k\) calls the bounded DFS visits every
non-cycle vertex exactly once (each is attached to a unique cycle vertex), so the
DFS batch is \(\bigO n\); building the doubled prefix and per-position values is
\(\bigO k\); the segment tree over \(2k\) leaves builds in \(\bigO k\) and the \(k\) window
queries cost \(\bigO{\log k}\) each.  With \(k\le n\le 2\cdot 10^{5}\) the total sits
well inside the 2\,s limit.  \textbf{Space} \(\bigO n\) adjacency plus \(\bigO k\) for the
cycle arrays and the \(8k\) segment tree; the \(\mathit{long\ long}\) values stay under
\(2\cdot 10^{14}\), and 256\,MB is never a risk.
\end{complexitybox}

\begin{hardeningbox}
\begin{itemize}[nosep,leftmargin=1.4em]
  \item \textbf{Everything is \code{long long}.}  With \(\ell_i\le 10^{9}\) summed along
        a path of up to \(n\) edges, the diameter reaches \(2\cdot 10^{14}\); a 32-bit
        int overflows silently exactly on the long-chain tests.
  \item \textbf{The \code{dfsAttach} recursion can depth \(n\).}  A 3-cycle with a
        \(2\cdot 10^{5}{-}3\) chain hanging off one vertex drives the fold that deep.
        Convert to an explicit stack (or raise the stack) before shipping; the two
        orderings are algebraically identical.
  \item \textbf{Leaves carry \(\mathrm{best}=-\infty\), not 0.}  A single-position
        interval has no pair \(a<b\).  Initialising it to 0 makes every window report
        at least 0 and can mask the real maximum when all \(\mathit{valMinus}\) are
        negative (common once \(S\) dominates \(h\)).
  \item \textbf{Query the doubled array, never wrap by hand.}  Writing
        \(\mathrm{dist}=\min(S_b-S_a,\ S_a-S_b{+}S_k)\) to avoid doubling answers a
        \emph{different} question -- it uses the shorter of the two arcs, i.e. the
        \emph{cycle's} distance, not the \emph{broken path's} distance.  The whole
        point of cutting edge \(w_j\) is that only one arc survives; the doubled
        prefix encodes exactly that, and hand-wrapping is dead end~3's bug.
  \item \textbf{Keep the \(\max(\mathit{dtreeBig},\cdot)\) fold.}  A bushy attached tree
        can make the answer an internal diameter no through-cycle window ever
        reaches; dropping the first argument is the construction in dead
        end~(3rd) and prints a too-small inconvenience.
\end{itemize}
\end{hardeningbox}

\begin{recallbox}[Chapter 29 in ten seconds]
\emph{A general graph is a tree plus redundant edges, and the craft is choosing
which redundancy to delete.  Bridges (\(\lemref{c31-lowchar}\)) are the edges you
cannot remove -- contract the rest to the bridge tree and Robbins
(\(\thmref{c31-robbins}\)) settles CF~118E; hinge vertices give the block-cut tree,
where the edge-stack pop (\(\thmref{c31-bccpop}\)) and the \(|E|{|}{=}|V|\) lone-cycle
lemma (\(\lemref{c31-onecycle}\)) answer CF~962F; an out-degree-one map peels to a
functional rho (\(\lemref{c31-peel}\)); edge-disjoint cycles contract to a cactus with
a ring transfer (\(\thmref{c31-cactustree}\), \(\lemref{c31-cycleorder}\)); and the four-test
signature (\(\defref{c31-signature}\), \(\defref{c31-pick}\)) decides which -- or whether to
leave for MST (\(\chref{c17}\)), counting (\(\chref{c32}\)) or dominators.  The Summit
CF~835F (\(\probref{c31-summit}\)) cuts the one cycle of a unicyclic graph and reads
the minimum tree diameter off a segment-tree window.}
\selftest{
  \item In a connected \(n\)-vertex \(n\)-edge graph, which edges are legal to delete
        while staying connected, and why? \emph{Exactly the unique cycle's edges -- every
        off-cycle edge is a bridge (\(\lemref{c31-lowchar}\)), and cycle rank \(|E|{-}|V|{+}1=1\)
        gives one cycle.}
  \item Why does the tree diameter after a cut split into \(\max(\max_i\mathit{dtree}_i,\)
        through-cycle)? \hfill\lv{5}
  \item The through-cycle term becomes \(\max_{a<b}(h_a-S_a)+(h_b+S_b)\) -- what data
        structure and what merge make that \(\bigO{\log k}\) per window? \emph{A segment tree
        over the doubled prefix with the monoid (mxMinus, mxPlus, best);
        \(\mathrm{best}=\max(\mathrm{L}.\mathrm{best},\mathrm{R}.\mathrm{best},\mathrm{L}.\mathrm{mxMinus}+\mathrm{R}.\mathrm{mxPlus})\) (\(\coderef{c31-summit}\)).}
  \item Greedy ``cut the longest cycle edge'' fails on which official sample, and by how
        much? \emph{Sample 2: greedy gives 24, the true answer is 18 -- off by 6.}
}
\end{recallbox}

\chaptersection{Tool Kit Inventory}

The machine assembled through Chapter~29.  Earlier chapters' rows are carried
forward; this chapter's are appended.  Pointers only.

\begin{center}
\small
\renewcommand{\arraystretch}{1.25}
\setlength{\tabcolsep}{3pt}
\begin{tabular}{@{}p{0.24\linewidth} l c p{0.22\linewidth} >{\raggedright\arraybackslash\scriptsize}p{0.17\linewidth}@{}}
\toprule
\textbf{Tool} & \textbf{Since} & \textbf{Lv} & \textbf{Trigger phrase} & \textbf{Where} \\
\midrule
\multicolumn{5}{l}{\itshape Carried forward} \\
Height, diameter by two sweeps & 12 & \lv{3} & ``deepest point / longest path in a tree'' & \chref{c09}, \thmref{c09-twosweep} \\
Bottom-up fold discipline & 13 & \lv{3} & ``children before parent, always'' & \chref{c10} \\
\(\tin/\tout\) DFS order & 5 & \lv{3} & ``is \(v\) inside my subtree'' & \chref{c24} \\
Segment tree range query & 19 & \lv{4} & ``max over an interval, associative merge'' & \chref{c13} \\
DSU / component merging & 17 & \lv{4} & ``contract, count components'' & \chref{c28} \\
\midrule
\multicolumn{5}{l}{\itshape Chapter 29 -- new this chapter} \\
Low-link bridge test & 29.1 & \lv{4} & ``must every path use this edge?'' & \lemref{c31-lowchar}, \coderef{c31-bridge} \\
Bridge tree contraction & 29.1 & \lv{4} & ``delete bridges, DP on components'' & \defref{c31-twiedge}, \figref{c31-bridgetree} \\
Robbins strong orientation & 29.1 & \lv{4} & ``orient so you can always return'' & \thmref{c31-robbins}, \probref{c31-roads} \\
Edge-stack block pop & 29.2 & \lv{5} & ``which blocks survive a hinge?'' & \thmref{c31-bccpop}, \coderef{c31-bcc} \\
Lone-cycle \(|E|{=}|V|\) test & 29.2 & \lv{5} & ``edge on exactly one cycle?'' & \lemref{c31-onecycle}, \probref{c31-cycles} \\
Functional rho peel & 29.3 & \lv{4} & ``one arrow per vertex, eventual cycle'' & \lemref{c31-peel}, \coderef{c31-peel} \\
Cactus contract + ring transfer & 29.4 & \lv{5} & ``each edge in at most one cycle'' & \thmref{c31-cactustree}, \lemref{c31-cycleorder} \\
Four-test signature, pick & 29.5 & \lv{5} & ``is there a tree worth extracting?'' & \defref{c31-signature}, \defref{c31-pick} \\
Window-monoid min diameter & 835F & \lv{6} & ``unicyclic, cut one edge, min diameter'' & \probref{c31-summit}, \coderef{c31-summit} \\
\bottomrule
\end{tabular}
\end{center}

\chaptersection{Related Topics}

\begin{itemize}[nosep, leftmargin=1.4em]
  \item \chref{c22} (Ch.~28, Tree Isomorphism) --- \textbf{backward, the assumption this
        chapter drops}: Ch.~28 assumed the input \emph{is} a tree and asked when two are
        the same shape; here the input is a general graph and the first job is to
        \emph{earn} a tree (bridges, blocks, cactus) before any Ch.~28 tool applies.
  \item \chref{c09} (Ch.~12, Height, Diameter, Centroid) --- \textbf{tool, used twice and
        never re-derived}: the two-sweep diameter is the black box the Summit runs on
        every attached tree, and the ``longest path through the root or inside a child''
        decomposition is exactly the \(\max(\mathit{dtree}_i,\cdot)\) split.
  \item \chref{c10} (Ch.~13, Tree DP) --- tool: once the tree-of-components is built
        (bridge tree, block-cut tree, cactus), every counting/optimisation answer is a
        bottom-up fold on it -- this chapter supplies the tree, Ch.~13 supplies the fold.
  \item \chref{c13} (Ch.~19, Segment Trees) --- tool in the Summit: the range-max window
        over the doubled prefix is an associative monoid on triples (mxMinus, mxPlus,
        best); no new machinery, just the right merge.
  \item \chref{c28} (Ch.~17, DSU / small-to-large) --- tool: component contraction (the
        bridge tree, the 920E complement components) is DSU work; DSU internals stay a
        black box cited back there.
  \item \chref{c30} (Ch.~26, Virtual Trees) --- contrast: a virtual tree \emph{compresses}
        marked vertices of an existing tree via LCA; this chapter \emph{decomposes} a
        general graph into a tree of components.  Both produce a small tree to DP on,
        from opposite starting assumptions.
  \item \chref{c17} (Ch.~31, Minimum Spanning Trees) --- \textbf{next-next neighbour, and
        the hand-off}: if the graph is weighted and you must \emph{choose} edges to
        minimise total cost, that is MST (Ch.~31), not a robustness decomposition; the
        signature's weight test sends you there (\(\defref{c31-signature}\)).
  \item \chref{c32} (Ch.~30, Counting Labelled Trees) --- \textbf{next chapter}: here one
        graph is given and you extract structure; there no graph is given at all and you
        \emph{count} how many trees meet a specification (Pr\"ufer, Matrix-Tree).
\end{itemize}

\begin{lastword}[Last Word --- Chapter 29]
\emph{``A graph is a tree that has been allowed to lie.  Redundant edges repeat
connections, hinge vertices hold pieces together, a single arrow folds the vertex
set back on itself, a cactus lets each edge break at most one rule.  Delete the
redundancy you can name -- the bridges, the articulation points, the in-degree
self-loops, the cycle blocks -- contract what remains, and you are handed back a
tree: the same tree Chapters~12 and~13 already taught you to run.''}

\medskip
Chapter~28 finished by insisting its input was a tree and asking only when two
such inputs were secretly identical.  This chapter refused that gift.  It opened
with the low-link scan (\(\lemref{c31-lowchar}\)) that finds the edges every path is
forced through, contracted them into the bridge tree, and let Robbins
(\(\thmref{c31-robbins}\)) decide CF~118E in one pass: orient inside a bridgeless
component, print 0 the instant a bridge survives.  It then traded edges for
vertices -- the edge-stack block pop (\(\thmref{c31-bccpop}\)) reads a block off the
stack at \(\mathrm{low}[v]\ge\tin[u]\), and CF~962F's whole question is which blocks
are lone cycles, answered by \(|E|{=}|V|\) (\(\lemref{c31-onecycle}\), brute-force
verified over \(7\cdot 10^{4}\) blocks before it entered print).  Section~29.3 turned
the arrow: a function on vertices is a rho of cycles with in-trees nailed to them,
peelable in \(\bigO n\) by in-degree (\(\lemref{c31-peel}\)).  Section~29.4 relaxed the
cactus and showed contracting every cycle leaves a tree (\(\thmref{c31-cactustree}\))
whose DP is a tree fold plus a ring transfer (\(\lemref{c31-cycleorder}\)).  Section~29.5
stated the four-test signature and the pick table so you can \emph{say no}: weights
to MST (\(\chref{c17}\)), enumeration to counting (\(\chref{c32}\)), directed flow cuts to
dominators -- and only one dominator preview, kept honest by admitting
Lengauer--Tarjan is not this chapter.

\medskip
The Summit CF~835F (\(\probref{c31-summit}\)) put it all together on a kingdom
\emph{almost} a tree: peel to the unique cycle, hand every cycle vertex to a Ch.~12
two-sweep for its height and internal diameter, break the cycle into a doubled
prefix, and let a Ch.~19 segment-tree monoid answer the best cut in
\(\bigO{k\log k}\).  The greedy `cut the longest cycle edge' fails on the official
sample by~6 -- a bridge-shaped lesson that this chapter exists to teach.

\medskip
Carry three reflexes forward:
\begin{enumerate}[nosep, leftmargin=1.6em]
  \item \emph{Name the redundancy before you touch it.}  Edge, vertex, arrow, cycle --
        the four tests of \(\defref{c31-signature}\) choose the decomposition, and
        choosing wrongly (bridge tree where blocks were needed, or vice-versa)
        loses the equality case and prints a wrong answer.
  \item \emph{Contract, then reuse.}  Once the redundant part is a node, the object
        you DP on is a tree -- Ch.~12 for heights, Ch.~13 for folds, Ch.~19 for windows.
        Nothing new is needed after the extraction; that is the entire point.
  \item \emph{Say `no' when the residual is not a tree.}  Two cycles sharing an edge
        break the ring transfer; a directed flow cut is not a bridge/block question;
        weighted minimisation is MST territory.  The hand-off table
        (\(\defref{c31-pick}\)) is as valuable as any algorithm here.
\end{enumerate}

\medskip
\textbf{Next:} \chref{c32} (Chapter~30, Counting Labelled Trees) removes the input
graph entirely.  You have spent this chapter earning a tree from a graph that was
given; the next asks how many trees meet a specification written on \(n\) labels --
Cayley's \(n^{n-2}\), Pr\"ufer's bijection between trees and length-\((n-2)\) sequences,
the Matrix--Tree theorem for spanning-tree counts of a general graph, and the
counter-intuition of `expected distance along a random walk'.  Where here the
tree was an answer, there the tree is a quantity; the tools shift from DFS to
linear-algebra, and the same Codeforces problem-set that taxed your bridge logic
here will tax your binomial coefficients next.
\end{lastword}
